% ============================================================
%  MTH636 – Game Theory | Lecture 24
%  Behavior Strategies, Mixed Strategies & Equivalence
% ============================================================
\documentclass[12pt, a4paper]{article}

% ============================================================
%  MTH636 -- Game Theory
%  Notes Preamble — Color-Coded, TikZ-Enhanced
% ============================================================

% --- Core packages ---
\usepackage[T1]{fontenc}
\usepackage[utf8]{inputenc}
\usepackage{lmodern}
\usepackage[margin=2.5cm, top=2.8cm, bottom=2.8cm]{geometry}
\usepackage{microtype}

% --- Math ---
\usepackage{amsmath, amssymb, amsthm, mathtools}
\usepackage{bm}        % bold math
\usepackage{cancel}    % strike-through in math

% --- Color & Boxes ---
\usepackage[dvipsnames, svgnames, x11names, table]{xcolor}
\usepackage[most]{tcolorbox}
\tcbuselibrary{breakable, skins, theorems}

% --- TikZ & Diagrams ---
\usepackage{tikz}
\usepackage{tikz-qtree}
\usetikzlibrary{arrows.meta, positioning, shapes.geometric, fit, backgrounds,
                decorations.pathreplacing, calc, matrix, trees}
\usepackage{pgfplots}
\pgfplotsset{compat=1.18}

% --- Tables ---
\usepackage{booktabs}
\usepackage{multirow}
\usepackage{array}
\usepackage{tabularx}
\usepackage{makecell}

% --- Typography & Layout ---
\usepackage{enumitem}
\usepackage{parskip}
\usepackage{fancyhdr}
\usepackage{titlesec}
\usepackage{titletoc}
\usepackage{hyperref}
\usepackage{graphicx}
\usepackage{float}
\usepackage{caption}

% --- Color Palette ---
\definecolor{defblue}{HTML}{1A5276}      % Definition header
\definecolor{defbluebg}{HTML}{D6EAF8}    % Definition background
\definecolor{thmgreen}{HTML}{1E8449}     % Theorem header
\definecolor{thmbg}{HTML}{D5F5E3}        % Theorem background
\definecolor{exorange}{HTML}{B7770D}     % Example header
\definecolor{exbg}{HTML}{FEF9E7}         % Example background
\definecolor{warnred}{HTML}{922B21}      % Warning/Key Insight header
\definecolor{warnbg}{HTML}{FADBD8}       % Warning background
\definecolor{noticegray}{HTML}{2C3E50}   % Note header
\definecolor{noticebg}{HTML}{EBF5FB}     % Note background
\definecolor{taketawayyellow}{HTML}{7D6608} % Key Takeaway header
\definecolor{takeawaybg}{HTML}{FDFEFE}   % Takeaway background
\definecolor{accentpurple}{HTML}{6C3483} % Accent / proof
\definecolor{lightgray}{HTML}{F2F3F4}    % shading

% ============================================================
%  TCOLORBOX ENVIRONMENTS
% ============================================================

% -- Definition (Blue) --
\newtcolorbox{defbox}[2][]{
  enhanced, breakable,
  colback=defbluebg, colframe=defblue, coltitle=white,
  fonttitle=\bfseries\sffamily,
  title={Definition: #2},
  attach boxed title to top left={yshift=-2mm, xshift=4mm},
  boxed title style={colback=defblue, rounded corners},
  arc=4pt, boxrule=1.2pt,
  left=6pt, right=6pt, top=8pt, bottom=6pt,
  #1
}

% -- Theorem / Proposition (Green) --
\newtcolorbox{thmbox}[2][]{
  enhanced, breakable,
  colback=thmbg, colframe=thmgreen, coltitle=white,
  fonttitle=\bfseries\sffamily,
  title={Theorem: #2},
  attach boxed title to top left={yshift=-2mm, xshift=4mm},
  boxed title style={colback=thmgreen, rounded corners},
  arc=4pt, boxrule=1.2pt,
  left=6pt, right=6pt, top=8pt, bottom=6pt,
  #1
}

% -- Proposition (Green variant) --
\newtcolorbox{propbox}[2][]{
  enhanced, breakable,
  colback=thmbg, colframe=thmgreen!70!black, coltitle=white,
  fonttitle=\bfseries\sffamily,
  title={Proposition: #2},
  attach boxed title to top left={yshift=-2mm, xshift=4mm},
  boxed title style={colback=thmgreen!70!black, rounded corners},
  arc=4pt, boxrule=1.2pt,
  left=6pt, right=6pt, top=8pt, bottom=6pt,
  #1
}

% -- Example (Orange/Yellow) --
\newtcolorbox{exbox}[2][]{
  enhanced, breakable,
  colback=exbg, colframe=exorange, coltitle=white,
  fonttitle=\bfseries\sffamily,
  title={Example: #2},
  attach boxed title to top left={yshift=-2mm, xshift=4mm},
  boxed title style={colback=exorange, rounded corners},
  arc=4pt, boxrule=1.2pt,
  left=6pt, right=6pt, top=8pt, bottom=6pt,
  #1
}

% -- Key Insight / Warning (Red) --
\newtcolorbox{warnbox}[2][]{
  enhanced, breakable,
  colback=warnbg, colframe=warnred, coltitle=white,
  fonttitle=\bfseries\sffamily,
  title={Key Insight: #2},
  attach boxed title to top left={yshift=-2mm, xshift=4mm},
  boxed title style={colback=warnred, rounded corners},
  arc=4pt, boxrule=1.2pt,
  left=6pt, right=6pt, top=8pt, bottom=6pt,
  #1
}

% -- Note (Gray-Blue) --
\newtcolorbox{notebox}[1][]{
  enhanced, breakable,
  colback=noticebg, colframe=noticegray!60, coltitle=white,
  fonttitle=\bfseries\sffamily,
  title={Note},
  attach boxed title to top left={yshift=-2mm, xshift=4mm},
  boxed title style={colback=noticegray!60, rounded corners},
  arc=4pt, boxrule=0.8pt,
  left=6pt, right=6pt, top=8pt, bottom=6pt,
  #1
}

% -- Key Takeaway --
\newtcolorbox{takebox}[1][]{
  enhanced, breakable,
  colback=takeawaybg, colframe=taketawayyellow, coltitle=white,
  fonttitle=\bfseries\sffamily,
  title={Key Takeaway},
  attach boxed title to top left={yshift=-2mm, xshift=4mm},
  boxed title style={colback=taketawayyellow, rounded corners},
  arc=4pt, boxrule=1.2pt,
  left=6pt, right=6pt, top=8pt, bottom=6pt,
  #1
}

% -- Intuition --
\newtcolorbox{intubox}[1][]{
  enhanced, breakable,
  colback=accentpurple!8, colframe=accentpurple!60, coltitle=white,
  fonttitle=\bfseries\sffamily,
  title={Intuition},
  attach boxed title to top left={yshift=-2mm, xshift=4mm},
  boxed title style={colback=accentpurple!80, rounded corners},
  arc=4pt, boxrule=0.8pt,
  left=6pt, right=6pt, top=8pt, bottom=6pt,
  #1
}

% ============================================================
%  PAYOFF MATRIX HELPER
%  Usage: inside tabular, use \payoff{a}{b} for (a,b)
% ============================================================
\newcommand{\payoff}[2]{$(#1,\;#2)$}
\newcommand{\nash}[2]{\underline{$#1$},\underline{$#2$}}  % underline NE payoffs
\newcommand{\Nset}{\mathbb{N}}
\newcommand{\Rset}{\mathbb{R}}
\newcommand{\Sset}{\mathcal{S}}
\newcommand{\Ui}{u_i}

% ============================================================
%  SECTION FORMATTING
% ============================================================
\titleformat{\section}
  {\large\bfseries\sffamily\color{defblue}}
  {\thesection}{1em}{}[\titlerule]

\titleformat{\subsection}
  {\normalsize\bfseries\sffamily\color{thmgreen!80!black}}
  {\thesubsection}{1em}{}

\titleformat{\subsubsection}
  {\normalsize\itshape\sffamily\color{exorange!80!black}}
  {\thesubsubsection}{1em}{}

% ============================================================
%  HEADER / FOOTER
% ============================================================
\pagestyle{fancy}
\fancyhf{}
\fancyhead[L]{\sffamily\small\color{defblue} MTH636 -- Game Theory}
\fancyhead[R]{\sffamily\small\color{defblue} \leftmark}
\fancyfoot[C]{\sffamily\small\thepage}
\renewcommand{\headrulewidth}{0.4pt}

% ============================================================
%  HYPERREF SETUP
% ============================================================
\hypersetup{
  colorlinks=true,
  linkcolor=defblue,
  urlcolor=thmgreen,
  citecolor=accentpurple,
  bookmarksnumbered=true,
  pdftitle={MTH636 Game Theory Notes},
  pdfauthor={MTH636 Student Notes}
}

% ============================================================
%  MISC HELPERS
% ============================================================
\newcommand{\BR}{\text{BR}}        % Best Response
\newcommand{\NE}{\text{NE}}        % Nash Equilibrium
\newcommand{\SPNE}{\text{SPNE}}    % Subgame Perfect NE
\newcommand{\Dom}{\text{dom}}      % Dominance
\newcommand{\mixs}[1]{\sigma_{#1}} % Mixed strategy
\newcommand{\pures}[1]{s_{#1}}     % Pure strategy


% ============================================================
\begin{document}

% Title
\begin{center}
  {\LARGE\bfseries\sffamily\color{defblue} MTH636 -- Game Theory}\\[6pt]
  {\Large\sffamily Lecture 24: Behavior Strategies, Mixed Strategies \& Equivalence}\\[4pt]
\end{center}
\vspace{0.5cm}
\hrule height 1.5pt
\vspace{0.3cm}

% ============================================================
\section{Complexity of Behavior vs.\ Mixed Strategies}

Recall that Player $i$ has $k$ information sets, and the $\ell$-th information set has $m_\ell$ many actions.

\begin{itemize}
  \item For a \textbf{behavior strategy} $b_i$, we have to specify
  \[
    \sum_{\ell=1}^{k}(m_\ell - 1) \quad \text{numbers.}
  \]

  \item For a \textbf{mixed strategy} $\sigma_i$, we have to specify
  \[
    \left(\prod_{\ell=1}^{k} m_\ell - 1\right) \quad \text{numbers.}
  \]
\end{itemize}

\begin{notebox}
  Behavior strategies are \emph{exponentially} more compact than mixed strategies, since
  $\sum_\ell (m_\ell-1) \ll \prod_\ell m_\ell - 1$ in general for large $k$.
\end{notebox}

% ============================================================
\section{Equivalence Definition}

\begin{defbox}{Equivalence of Mixed and Behavior Strategies}
  A mixed strategy $\sigma_i$ and a behavior strategy $b_i$ of Player $i$ in an extensive form game $\Gamma$ are said to be \textbf{equivalent} if, for every mixed/behavior strategy vector $\sigma_{-i}$ of the other players and for all vertices $n$ of the game tree,
  \[
    \rho\bigl(n;\,(\sigma_i,\sigma_{-i})\bigr) = \rho\bigl(n;\,(b_i,\sigma_{-i})\bigr),
  \]
  where $\rho(n;\cdot)$ denotes the probability of reaching vertex $n$ under the given strategy profile.
\end{defbox}

\subsection{Illustrative Game Tree}

Consider the following three-player extensive form game tree (Player I at root $n_0$, Player II at $n_1$, and Player I again at the information set $\{n_3,n_4\}$):

\medskip
\noindent
\textbf{Tree structure:}
\begin{itemize}
  \item $n_0$ (Player I): actions $A$ (to $n_1$) and $B$ (to $n_2$).
  \item $n_1$ (Player II): actions $L$ (to $n_3$) and $R$ (to $n_4$). Nodes $n_3$ and $n_4$ belong to the same information set of Player I (dashed box).
  \item At the information set $\{n_3, n_4\}$ (Player I): actions $\ell$ (to $n_5$ or $n_7$) and $r$ (to $n_6$ or $n_8$).
  \item Terminal nodes: $n_2, n_5, n_6, n_7, n_8$.
\end{itemize}

% ============================================================
\section{Remarks on Equivalence}

\subsection*{Remark R-1}

Suppose $b_i$ (a behavior strategy) and $\sigma_i$ (a mixed strategy) are equivalent. Then, for all mixed/behavior strategy vectors $\sigma_{-i}$ of the other players $\neq i$,
\[
  \textcolor{red}{U_j(\sigma_i,\sigma_{-i}) = U_j(b_i,\sigma_{-i})}.
\]
More explicitly, this gives:
\[
  U_j\!\left(\sigma_i,\sigma_j^{\vphantom{x}},\sigma_{-\{i,j\}}\right)
  = U_j\!\left(b_i,\sigma_j^{\vphantom{x}},\sigma_{-\{i,j\}}\right).
\]

\subsection*{Remark R-2}

Let $\sigma = (\sigma_i)_{i\in N}$ be a mixed strategy profile and $b = (b_i)_{i\in N}$ be a behavior strategy profile. Suppose $b_i$ is equivalent to $\sigma_i$ for all $i$. Then we have
\[
  U_i(\sigma) = U_i(b) \quad \forall\, i.
\]

% ============================================================
\section{Worked Example: Verifying Equivalence}

\subsection{Setup}

Consider the game tree with:
\begin{itemize}
  \item $n_0$ (Player I, root): actions $A$ (prob.\ $\tfrac{1}{2}$) and $B$ (prob.\ $\tfrac{1}{2}$) under mixed strategy $\sigma_1$.
  \item $n_1$ (Player II): actions $L$ and $R$.
  \item Information set $\{n_3, n_4\}$ (Player I): actions $\ell$ and $r$.
  \item Behavior strategy of Player II: $\sigma_2(L) = 1$.
\end{itemize}

Under the mixed strategy $\sigma_1$ (Player I plays $A$ with prob.\ $\tfrac{1}{2}$ and $B$ with prob.\ $\tfrac{1}{2}$; within $B$, plays $\ell$ with prob.\ $\tfrac{1}{3}$ and $r$ with prob.\ $\tfrac{2}{3}$):
\begin{align*}
  \sigma_1(A\ell) &= \tfrac{1}{2} \times \tfrac{1}{3}, &
  \sigma_1(Ar)    &= \tfrac{1}{2} \times \tfrac{2}{3}, \\
  \sigma_1(B\ell) &= \tfrac{1}{2} \times \tfrac{1}{3}, &
  \sigma_1(Br)    &= \tfrac{1}{2} \times \tfrac{2}{3}.
\end{align*}

The reaching probabilities must satisfy:
\[
  \rho\bigl(n_2;\,(\sigma_1,\sigma_2)\bigr) = \rho\bigl(n_2;\,(b_1,\sigma_2)\bigr),
\]
\[
  \rho\bigl(n_i;\,(\sigma_1,\sigma_2)\bigr) = \rho\bigl(n_i;\,(b_1,\sigma_2)\bigr)
  \quad \forall\, n_i \in L,
\]
where $L$ denotes the set of all leaves.

\subsection{Constructing an Equivalent Mixed Strategy from a Behavior Strategy}

Given a behavior strategy $b_i$, we can construct an equivalent mixed strategy $\sigma_i$ via:
\[
  \sigma_i(s_i) = \prod_{h_i \in H_i} b_i\!\bigl(s_i(h_i),\, h_i\bigr),
\]
i.e., the probability of a pure strategy $s_i$ under $\sigma_i$ is the product of the behavior strategy probabilities at each information set.

\medskip
\noindent\textbf{Example.} Using the tree above with actions $A$, $B$, $\ell$, $r$ and probability labels:
\begin{itemize}
  \item Probabilities at $n_5, n_6$: $\ell \mapsto \tfrac{1}{3}$, $r \mapsto \tfrac{2}{3}$.
  \item Probabilities at $n_7, n_8$: $\ell \mapsto \tfrac{1}{2}$, $r \mapsto \tfrac{2}{3}$.
\end{itemize}

Computing reaching probabilities:
\[
  \sigma_1(A\ell) = 0,\quad \sigma_1(Ar) = 0, \quad
  \sigma_1(B\ell) = \tfrac{1}{2} \cdot \alpha, \quad \sigma_1(Br) = \tfrac{1}{2}(1-\alpha),
\]
for some free parameter $\alpha \in [0,1]$.

\begin{notebox}
  This equivalence is \textbf{not unique} in the sense that given a behavior strategy, there can be \textbf{multiple} mixed strategies which are equivalent to it. The parameter $\alpha$ above is free, illustrating this non-uniqueness.
\end{notebox}

% ============================================================
\section{Theorem: Equivalence Always Exists (Behavior $\Rightarrow$ Mixed)}

\begin{thmbox}{Existence of Equivalent Mixed Strategy}
  Given a behavior strategy $b_i$, there always exists an equivalent mixed strategy $\sigma_i$.
\end{thmbox}

\begin{proof}
  Given $b_i$, construct $\sigma_i$ by setting
  \[
    \sigma_i(s_i) = \prod_{h_i \in H_i} b_i\!\bigl(s_i(h_i),\, h_i\bigr).
  \]
  One then verifies that this $\sigma_i$ produces the same reaching probability at every vertex of the game tree as $b_i$ does.
\end{proof}

% ============================================================
\section{Mixed $\Rightarrow$ Behavior: The Non-Trivial Direction}

The non-trivial question is: given a mixed strategy $\sigma_i$, does an equivalent behavior strategy exist?

\subsection{Example}

Consider a game tree where Player I has two information sets $h_1 = \{n_0\}$ (root) and $h_2 = \{n_1, n_2, n_3\}$ (a three-node information set). Player I has pure strategies $A\ell$, $Ar$, $B\ell$, $Br$, and the mixed strategy satisfies:
\begin{align*}
  \sigma_1(A\ell) &= \tfrac{1}{2}, & \sigma_1(B\ell) &= \tfrac{1}{2}, \\
  \sigma_1(Ar)    &= 0,            & \sigma_1(Br)    &= 0.
\end{align*}

Let $b_1(A;\, n_0) = u$ (probability of action $A$ at root) and $b_1(\ell;\, h_2) = y$ (probability of action $\ell$ at $h_2$). Then:
\begin{align*}
  \rho(n_3;\,\sigma_1) &= \tfrac{1}{2} = uy, & 0 &= (1-u)y, \\
  \rho(n_4;\,\sigma_1) &= 0 = u(1-y),        & \tfrac{1}{2} &= (1-u)(1-y).
\end{align*}

From the system $uy = \tfrac{1}{2}$ and $(1-u)y = 0$: we get $y = 0$ (from the second equation if $u \neq 1$) but $uy = \tfrac{1}{2}$ requires $y > 0$. This contradiction shows the system has \textbf{no solution}, demonstrating that:

\begin{warnbox}{No Equivalent Behavior Strategy Without Perfect Recall}
  If we do \textbf{not} have \emph{perfect recall}, then there exists a mixed strategy which does \textbf{not} have an equivalent behavior strategy.
\end{warnbox}

\end{document}
